% -----------------------------------------------------------------------------
% This file is automatically generated.
% Do not edit manually.
% Source: notebooks/support/hypothesis-testing/support.Rmd
% -----------------------------------------------------------------------------


\noindent This result is obtained using the default
Hypergeometric distribution. To obtain the sample
size with the binomial approximation, as in Table
5.3, we specify the distribution when setting up
the attribute object.

\par\addvspace{\topsep}
\begingroup
\fvset{listparameters={%
  \setlength{\topsep}{0pt}%
  \setlength{\partopsep}{0pt}%
  \setlength{\parsep}{0pt}%
  \setlength{\itemsep}{0pt}%
}}
\begin{Verbatim}[breaklines=true,formatcom=\color{ada_blue}]
mySample2 = att_sample(alpha=0.1, tdr=0.05, c=2, 
dist="binom").size()
mySample2.n
\end{Verbatim}
\begin{Verbatim}[breaklines=true,formatcom=\color{ada_light_blue}]
[1] 105
\end{Verbatim}
\endgroup
\par\addvspace{\topsep}
